\section{Introduction}
\label{sec:introduction}

Mass spectrometry imaging (MSI) records a full mass spectrum at each sampled
location of a tissue section, so the spatial distribution of many molecules can
be studied at once without choosing targets in
advance~\citep{caprioli1997maldi,schwamborn2010molecular,buchberger2018mass}.
The data are large and high-dimensional. The collections used in this report
contain between about a thousand and several thousand measured locations per
section, and between 1{,}481 and 85{,}062 mass-to-charge ($m/z$) bins per
spectrum. Most analyses therefore begin by selecting a small set of informative
peaks. Classical peak picking applies smoothing, normalisation and
signal-to-noise thresholds to each spectrum~\citep{gibb2012maldiquant,deininger2011normalization},
so it ignores \emph{where} in the tissue a signal occurs.

This report follows the distinction drawn in the proposal between
\emph{peak picking}, which applies rule-based thresholds to each spectrum, and
\emph{peak learning}, which treats peak selection as an unsupervised
representation-learning problem~\citep{li2022machinelearningpeakpicking,abdelmoula2021peak}.
Peak picking depends on manually set thresholds and can discard weak but
informative ions~\citep{lieb2020sparse}. Peak learning instead fits a model to
the spectra and ranks $m/z$ bins by what the model uses. msiPL trains a variational autoencoder (VAE)
on individual spectra and reads peaks from the learned
weights~\citep{abdelmoula2021peak}. S3PL adds spatial information by convolving
over small tissue patches and reconstructing the whole
neighbourhood~\citep{weigand2026spatial}. That raises two questions. First, can
cheaper local spatial context, such as the mean spectrum of a $3\times3$
neighbourhood, give a lightweight VAE a similar benefit? This was the original
aim of this project. The proposal, titled \emph{Representation Learning for
Autonomous Peak Detection in Mass Spectrometry Imaging}, set out to build
\emph{Spatial-msiPL}: an msiPL-style VAE with neighbourhood input and a
spatial-coherence penalty, intended as a computationally efficient,
spatially aware alternative to 3D convolutional networks. Second, once a model has been fitted, which part of the
pipeline actually determines peak quality: the representation, the
segmentation built on it, or the attribution method used to rank peaks?

This report addresses both questions experimentally. We built Spatial-msiPL
and its variants: VAEs
that encode either a single spectrum or the spectrum together with its
neighbourhood (uniform averaging, learned attention, larger windows, masked
centres). We segmented each VAE's latent space with a Gaussian mixture model
(GMM) and ranked peaks with Integrated Gradients (IG)~\citep{sundararajan2017axiomatic}
of the GMM posteriors. The evaluation was designed to be fair rather than
favourable:
\begin{itemize}
  \item every method selected the same number of bins per section;
  \item decision rules were written down before the corresponding results existed;
  \item reproduced msiPL and S3PL baselines and simple model-free rankings
        (mean intensity, variance, Moran's~$I$) were scored the same way.
\end{itemize}
The design turned out to matter. Holding the segmentation fixed made it
possible to separate the value of the learned segmentation from the value of
the attribution method.

The main findings are as follows.
\begin{itemize}
  \item \textbf{Ranking.} Correlating each ion image directly with the GMM
        posterior maps outperformed IG in 136 of 140 paired comparisons across
        three training seeds. This difference was statistically supported in all
        six dataset--model comparisons.
  \item \textbf{Spatial information.} Neighbourhood input to the encoder gave
        no consistent, spatially specific gain. Spatially aware ranking
        statistics did: correlation with the segmentation maps was strongest on
        glioblastoma, and Moran's~$I$ was strongest on renal cell carcinoma.
        This is the sense in which peak ranking benefits from being spatially
        aware.
  \item \textbf{Encoder.} A linear five-component PCA segmentation performed
        comparably to the VAE, so the nonlinear encoder was not shown to be
        necessary.
  \item \textbf{Reproducibility.} The published S3PL glioblastoma score lay
        outside the distribution of ten reproduced seeds. Seeds changed S3PL
        mainly through weight initialisation, which switched individual sections
        between distinct solutions.
\end{itemize}
All conclusions are bounded to this pipeline, these datasets and the
mask-agreement metric used by the field (\cref{sec:discussion}).

The contributions of this report are:
\begin{enumerate}
  \item a predeclared, matched-budget evaluation protocol for learned MSI peak
        ranking, applied to three datasets with reproduced baselines and
        model-free references;
  \item a ranking ablation showing that segment-then-correlate outperforms
        gradient attribution in this pipeline, reported with paired statistics;
  \item controlled tests of spatial context at the encoder and segmentation
        stages that bound where spatial information helps;
  \item a statistical account of the S3PL reproduction gap and its seed
        sensitivity.
\end{enumerate}

\Cref{sec:background} reviews MSI peak selection and the methods compared.
\Cref{sec:questions} states the research questions and how they changed from
the proposal. \Cref{sec:methodology} describes the data, models, evaluation,
statistics and computation. \Cref{sec:results} presents the results.
\Cref{sec:discussion} interprets them and discusses limitations, and
\cref{sec:conclusion} concludes.
